\begin{center}\begin{tikzpicture}[font=\small,>=stealth,line join=round]
\draw[gray!40,rounded corners=3pt] (0,0) rectangle (14.2,8.3);\node at (7.1,7.700000000000001) {连续曲线：从正高度到负高度，途中必碰横轴};
\begin{scope}[shift={(1.7666666666666666,3.8842857142857143)},x=2.2222222222222223cm,y=1.3142857142857143cm]\draw[->,gray] (-0.3,0)--(2.4,0) node[right]{$x$};\draw[->,gray] (0,-1.7)--(0,1.8) node[above]{$y$};\draw[blue,thick,solid,domain=-0.1:2.1,samples=180] plot(\x,{cos(deg(pi*\x))});\fill[orange!85!black] (0,1) circle(2.2pt);\fill[orange!85!black] (1,-1) circle(2.2pt);\fill[orange!85!black] (0.5,0) circle(2.2pt);\end{scope}\node[anchor=west,align=left,text width=5.0cm] at (8.4,4.0) {$f=\cos\pi x,\ a=1$\\[9pt]$f(0)=1,\ f(1)=-1$\\[9pt]在 $[0,1]$ 连续\\[9pt]$f(\frac12)=0$\\[9pt]至少有一个零点};
\node[text width=12.6cm,align=center] at (7.1,.65) {连续是保证中途取到0的关键条件；端点本来为0时也算碰到横轴。};
\end{tikzpicture}\end{center}
